\documentclass[10pt]{beamer}

\usepackage{appendixnumberbeamer}
\usepackage{booktabs}
\usepackage[scale=2]{ccicons}
\usepackage{pgfplots}
\usepackage{xspace}
\usepackage{bookmark}
\usepackage{amssymb}
\usepackage{mathtools}
\usepackage[normalem]{ulem}
\usepackage[T1]{fontenc}
\usepackage[sfdefault,book]{FiraSans}
\usepackage{FiraMono}
\usepackage{fontawesome}
\usepackage{hyperref}
\usepackage{listings}

\renewcommand{\thefootnote}{\fnsymbol{footnote}}
\renewcommand{\MintedPygmentize}{/home/thiago/.local/bin/pygmentize}

\definecolor{light-gray}{gray}{0.95}
\renewcommand\lstlistingname{Código}
\DeclareCaptionFormat{listing} {
  \parbox{\textwidth}{\hspace{-0.2cm}#1#2#3}
}
\DeclareCaptionFont{black}{\color{black}}
\captionsetup[lstlisting]{
  format=listing,
  labelfont=black,
  textfont=black,
  singlelinecheck=true,
  margin=0pt,
  font={tt,footnotesize,bf}
}

\lstset{
  basicstyle=\footnotesize\ttfamily,
  escapeinside={\%*}{*)},
  mathescape=true,
  showspaces=false,
  showtabs=false,
  showstringspaces=false,%
  rulesepcolor=\color{black},
  frame=shadowbox,
  upquote=true,
  literate=
  {á}{{\'a}}1 {é}{{\'e}}1 {í}{{\'i}}1 {ó}{{\'o}}1 {ú}{{\'u}}1
  {Á}{{\'A}}1 {É}{{\'E}}1 {Í}{{\'I}}1 {Ó}{{\'O}}1 {Ú}{{\'U}}1
  {à}{{\`a}}1 {è}{{\`e}}1 {ì}{{\`i}}1 {ò}{{\`o}}1 {ù}{{\`u}}1
  {À}{{\`A}}1 {È}{{\'E}}1 {Ì}{{\`I}}1 {Ò}{{\`O}}1 {Ù}{{\`U}}1
  {ä}{{\"a}}1 {ë}{{\"e}}1 {ï}{{\"i}}1 {ö}{{\"o}}1 {ü}{{\"u}}1
  {Ä}{{\"A}}1 {Ë}{{\"E}}1 {Ï}{{\"I}}1 {Ö}{{\"O}}1 {Ü}{{\"U}}1
  {â}{{\^a}}1 {ê}{{\^e}}1 {î}{{\^i}}1 {ô}{{\^o}}1 {û}{{\^u}}1
  {Â}{{\^A}}1 {Ê}{{\^E}}1 {Î}{{\^I}}1 {Ô}{{\^O}}1 {Û}{{\^U}}1
  {Ã}{{\~A}}1 {ã}{{\~a}}1 {Õ}{{\~O}}1 {õ}{{\~o}}1
  {œ}{{\oe}}1 {Œ}{{\OE}}1 {æ}{{\ae}}1 {Æ}{{\AE}}1 {ß}{{\ss}}1
  {ű}{{\H{u}}}1 {Ű}{{\H{U}}}1 {ő}{{\H{o}}}1 {Ő}{{\H{O}}}1
  {ç}{{\c c}}1 {Ç}{{\c C}}1 {ø}{{\o}}1 {å}{{\r a}}1 {Å}{{\r A}}1
  {€}{{\euro}}1 {£}{{\pounds}}1 {«}{{\guillemotleft}}1
  {»}{{\guillemotright}}1 {ñ}{{\~n}}1 {Ñ}{{\~N}}1 {¿}{{?`}}1
}

\setminted{
  fontsize=\footnotesize,
  style=bw,
  frame=single,
  labelposition=topline,
}

\pretitle{\begin{center}\normalsize\bfseries\MakeUppercase{Programação 2 -- }\MakeUppercase}
\author{\normalsize Prof. Thiago Cavalcante}
\date{\vspace{-4ex}}

\geometry{
  top=0.0cm,
  left=1.0cm,
  right=1.0cm,
  bottom=1.0cm
}

\setlength\columnsep{30pt}

\usetikzlibrary{arrows}

\makeatletter
\AddEnumerateCounter{\PaddingUp}{\two@digits}{A00}
\AddEnumerateCounter{\PaddingDown}{\two@digits}{A00}
\newcommand\PaddingUp[1]{\expandafter\two@digits\csname c@#1\endcsname}
\newcommand\PaddingDown[1]{\PaddingUp{#1}\addtocounter{#1}{-2}}
\makeatother

\makeatletter
\def\enumalphalphcnt#1{\expandafter\@enumalphalphcnt\csname c@#1\endcsname}
\def\@enumalphalphcnt#1{\AlphAlph{#1}}
\makeatother
\AddEnumerateCounter{\enumalphalphcnt}{\@enumalphalphcnt}{aa}

\pagenumbering{gobble}

\title{Programação 2}
\author{Thiago Cavalcante  -- thiago.cavalcante@penedo.ufal.br}
\institute{Universidade Federal de Alagoas -- UFAL \\ Campus Arapiraca \\ Unidade de Ensino de Penedo}
\titlegraphic{\hfill\includegraphics[height=1.5cm]{images/brasao-ufal.eps}}


\subtitle{Aula 6}
\date{22 de janeiro de 2020}

\begin{document}

\maketitle

\begin{frame}{Arquivos}
  \LARGE
  \textbf{Arquivos são coleções de bytes} armazenados em um dispositivo de armazenamento (como um disco rídigo, por exemplo)
  \vfill
  Um arquivo pode ser interpretado como um texto, uma imagem, um vídeo etc.
\end{frame}

\begin{frame}{Arquivos}
  \LARGE
  \textbf{Arquivos são coleções de bytes} armazenados em um dispositivo de armazenamento (como um disco rídigo, por exemplo)
  \vfill
  Um arquivo pode ser interpretado como um texto, uma imagem, um vídeo etc.
\end{frame}

\begin{frame}{Arquivos}
  \huge
  \textbf{Vantagens}
  \vfill
  \Large
  \begin{itemize}
    \item Armazenamento durável
    \item Grandes quantidades de informação
    \item Acesso sequencial ou não
    \item Acesso concorrente (vários programas)
  \end{itemize}
\end{frame}

\begin{frame}{Arquivos}
  \LARGE
  \centering
  \textbf{arquivos texto $\times$ arquivos binários}
\end{frame}

\begin{frame}{Arquivos}
  \huge
  \textbf{Arquivos texto}
  \vfill
  \Large
  \begin{itemize}
    \item Podem ser mostrados \textbf{diretamente na tela} ou modificados em um editor de texto como o \textbf{bloco de notas}
    \item Dados são gravados como \textbf{caracteres de 8 bits da tabela ASCII} (é feita uma conversão)
  \end{itemize}
\end{frame}

\begin{frame}{Arquivos}
  \huge
  \textbf{Arquivos binários}
  \vfill
  \Large
  \begin{itemize}
    \item Armazenam uma \textbf{sequência de bits} sujeita às convenções dos programas que a geraram
    \item São gravados \textbf{exatamente como estão organizados na memória}, sem conversão
    \item Exemplos: arquivos executáveis, arquivos compactados etc.
  \end{itemize}
\end{frame}

\begin{frame}{Arquivos}
  \huge
  \textbf{Leitura e escrita}
  \vfill
  \LARGE
  Funções da biblioteca \textbf{\texttt{stdio.h}}
  \vfill
  \begin{itemize}
    \item Abrir/fechar
    \item Leitura/escrita de caracteres/bytes
  \end{itemize}
\end{frame}

\begin{frame}[fragile]{Arquivos}
  \huge
  \textbf{Ponteiro para arquivo}
  \vfill
  \begin{lstlisting}[language=C]
FILE *p;
  \end{lstlisting}
  \vfill
  \large
  Aponta para uma área na memória chamada de \textbf{buffer}, que contém informações sobre o arquivo aberto
\end{frame}

\begin{frame}[fragile]{Arquivos}
  \huge
  \textbf{Abrindo um arquivo: \texttt{fopen}}
  \vfill
  \large
  \begin{lstlisting}[language=C]
FILE *fopen(char *nome_do_arquivo, char *modo)
  \end{lstlisting}
  \vfill
  \setlength{\leftmargini}{0pt}
  \begin{itemize}
      \item [] \textbf{\texttt{nome\_do\_arquivo}} pode ser um caminho \textbf{absoluto} ou \textbf{relativo}, com relação à localização do programa
      \item [] Retorna \textbf{\texttt{NULL}} caso ocorra um erro
  \end{itemize}
\end{frame}

\begin{frame}{Arquivos}
  \large
  \textbf{Localização do programa}

  \texttt{"C:\textbackslash Programas\textbackslash programa.c"}
  \vfill
  \textbf{Localização do arquivo (absoluta)}

  \texttt{"C:\textbackslash Programas\textbackslash Arquivos\textbackslash arquivo.txt"}
  \vfill
  \textbf{Localização do arquivo (relativa)}

  \texttt{".\textbackslash Arquivos\textbackslash arquivo.txt"}
\end{frame}

\begin{frame}{Arquivos}
  \footnotesize
  \begin{center}
    \begin{tabular}{@{}ll@{}}
      \toprule
      \textbf{Modo} & \textbf{Função} \\
      \midrule
      \texttt{"r"} & Leitura. Arquivo deve existir. \\
      \texttt{"w"} & Escrita. Cria arquivo, se não houver. Apaga o anterior, se ele existir. \\
      \texttt{"a"} & Escrita. Os dados serão adicionados no fim do arquivo (\textit{append}). \\
      \texttt{"r+"} & Leitura/escrita. O arquivo deve existir e pode ser modificado. \\
      \texttt{"w+"} & Leitura/escrita. Cria arquivo se não houver. Apaga o anterior se ele existir. \\
      \texttt{"a+"} & Leitura/escrita. Os dados serão adicionados no fim do arquivo (\textit{append}). \\
      \bottomrule
    \end{tabular}
  \end{center}
  \vfill
  \large
  Para arquivos binários, adiciona-se um \textbf{\texttt{b}} ao modo
\end{frame}

\begin{frame}[fragile]{Arquivos}
  \huge
  \textbf{Saindo de um programa em caso de erro}
  \vfill
  \Large
  \begin{lstlisting}[language=C]
void exit(int codigo_de_retorno)
  \end{lstlisting}
  \vfill
  \large
  Convenção: \texttt{\textbf{codigo\_de\_retorno}} é \texttt{\textbf{0}} para um término normal e outro número caso ocorra um problema
\end{frame}

\begin{frame}[fragile]{Arquivos}
  \huge
  \textbf{Fechando um arquivo: \texttt{fclose}}
  \vfill
  \large
  \begin{lstlisting}[language=C]
int fclose(FILE *fp)
  \end{lstlisting}
  \vfill
  \setlength{\leftmargini}{0pt}
  \begin{itemize}
      \item [] Retorna \texttt{\textbf{0}} no caso de sucesso e outro número, caso contrário
      \item [] Ao fechar um arquivo, todo caractere que tenha permanecido no buffer é \textbf{gravado} (também ocorre quando o buffer fica cheio)
      \item [] A função \texttt{\textbf{exit}} fecha todos os arquivos abertos
  \end{itemize}
\end{frame}

\begin{frame}[fragile]{Arquivos}
  \huge
  \textbf{Escrevendo um char: \texttt{fputc}}
  \vfill
  \large
  \begin{lstlisting}[language=C]
int fputc(int c, FILE *fp)
  \end{lstlisting}
  \vfill
  \setlength{\leftmargini}{0pt}
  \begin{itemize}
      \item [] Retorna a constante \textbf{\texttt{EOF}}, se houver erro na escrita, ou o \textbf{próprio caractere}, se ele foi escrito com sucesso
      \item [] Após a escrita, o \textbf{indicador de posição interna} é avançado em um caractere (pode ser visto como o \textbf{cursor})
      \item [] Pode-se escrever na tela com o arquivo \textbf{\texttt{stdout}}
  \end{itemize}
\end{frame}

\begin{frame}[fragile]{Arquivos}
  \huge
  \textbf{Lendo um char: \texttt{fgetc}}
  \vfill
  \large
  \begin{lstlisting}[language=C]
int fgetc(FILE *fp)
  \end{lstlisting}
  \vfill
  \setlength{\leftmargini}{0pt}
  \begin{itemize}
    \item [] Retorna a constante \textbf{\texttt{EOF}}, se houver erro na leitura, ou o \textbf{próprio caractere}, se ele foi lido com sucesso
    \item [] Após a leitura, o \textbf{indicador de posição interna} é avançado em um caractere
      \item [] Pode-se ler do teclado com o arquivo \textbf{\texttt{stdin}}
    \end{itemize}
\end{frame}

\begin{frame}[fragile]{Arquivos}
  \huge
  \textbf{Fim do arquivo: \texttt{feof}}
  \vfill
  \large
  \begin{lstlisting}[language=C]
int feof(FILE *fp)
  \end{lstlisting}
  \vfill
  \setlength{\leftmargini}{0pt}
  \begin{itemize}
      \item [] Retorna \texttt{\textbf{0}} caso o fim ainda não tenha sido atingido e outro número, caso contrário
  \end{itemize}
\end{frame}

\begin{frame}[fragile]{Arquivos}
  \huge
  \textbf{Posição atual: \texttt{ftell}}
  \vfill
  \large
  \begin{lstlisting}[language=C]
long int ftell(FILE *fp)
  \end{lstlisting}
  \vfill
  \setlength{\leftmargini}{0pt}
  \begin{itemize}
      \item [] Para \textbf{arquivos binários}, o valor retornado indica o \textbf{número de bytes lidos} a partir do início do arquivo
      \item [] Para arquivos \textbf{texto}, \textbf{não existe garantia} de que o valor retornado seja o número exato de bytes lidos a partir do início do arquivo (mas esse valor pode ser usado para se retornar ao mesmo ponto, de qualquer forma)
      \item [] Em caso de erro, retorna -1
  \end{itemize}
\end{frame}

\begin{frame}[fragile]{Arquivos}
  \huge
  \textbf{Escrevendo uma string: \texttt{fputs}}
  \vfill
  \large
  \begin{lstlisting}[language=C]
int fputs(char *str,FILE *fp)
  \end{lstlisting}
  \vfill
  \setlength{\leftmargini}{0pt}
  \begin{itemize}
      \item [] Retorna a constante \textbf{\texttt{EOF}}, se houver erro na escrita, ou um valor diferente de \texttt{\textbf{0}}, caso contrário
      \item [] Não insere uma nova linha \texttt{`\textbackslash n'} (responsabilidade do programador)
  \end{itemize}
\end{frame}

\begin{frame}[fragile]{Arquivos}
  \huge
  \textbf{Lendo uma string: \texttt{fgets}}
  \vfill
  \large
  \begin{lstlisting}[language=C]
char *fgets(char *str, int tamanho, FILE *fp)
  \end{lstlisting}
  \vfill
  \setlength{\leftmargini}{0pt}
  \begin{itemize}
      \item [] Retorna \textbf{\texttt{NULL}}, se houver erro na leitura, ou o ponteiro para o primeiro caractere da string, caso contrário
      \item [] Lê até encontrar uma nova linha \texttt{`\textbackslash n'} (  que fará parte da string) ou até receber \texttt{\textbf{tamanho - 1}} caracteres (o último é sempre \texttt{`\textbackslash 0'})
  \end{itemize}
\end{frame}

\begin{frame}[fragile]{Arquivos}
  \huge
  \textbf{Escrevendo blocos de bytes: \texttt{fwrite}}
  \vfill
  \large
  \begin{lstlisting}[language=C]
int fwrite(void *buffer, int nro_de_bytes,
           int count, FILE *fp)
  \end{lstlisting}
  \vfill
  \setlength{\leftmargini}{0pt}
  \begin{itemize}
      \item [] Retorna o número de unidades escritas com sucesso (pode ser menor que \textbf{\texttt{count}} \faExclamationTriangle)
  \end{itemize}
\end{frame}

\begin{frame}[fragile]{Arquivos}
  \huge
  \textbf{Lendo blocos de bytes: \texttt{fread}}
  \vfill
  \large
  \begin{lstlisting}[language=C]
int fread(void *buffer, int nro_de_bytes,
          int count, FILE *fp)
  \end{lstlisting}
  \vfill
  \setlength{\leftmargini}{0pt}
  \begin{itemize}
      \item [] Retorna o número de unidades escritas com sucesso (pode ser menor que \textbf{\texttt{count}} \faExclamationTriangle)
  \end{itemize}
\end{frame}

\begin{frame}{Arquivos}
  \huge
  \textbf{Boa prática: escrever no arquivo o tamanho da string/array antes do conteúdo (facilita na hora de ler, depois)}
\end{frame}

\begin{frame}[fragile]{Arquivos}
  \huge
  \textbf{Escrevendo dados formatados: \texttt{fprintf}}
  \vfill
  \large
  \begin{lstlisting}[language=C]
int fprintf(FILE *fp, "tipos de saída",
            lista de variáveis)
  \end{lstlisting}
  \vfill
  \setlength{\leftmargini}{0pt}
  \begin{itemize}
      \item [] Retorna o total de caracteres escritos, em caso de sucesso, ou um número negativo, caso contrário
  \end{itemize}
\end{frame}

\begin{frame}[fragile]{Arquivos}
  \huge
  \textbf{Lendo dados formatados: \texttt{fscanf}}
  \vfill
  \large
  \begin{lstlisting}[language=C]
int fscanf(FILE *fp, "tipos de entrada",
           lista de variáveis)
  \end{lstlisting}
  \vfill
  \setlength{\leftmargini}{0pt}
  \begin{itemize}
      \item [] Retorna o total de caracteres lidos, em caso de sucesso, ou a constante \textbf{\texttt{EOF}}, caso contrário
  \end{itemize}
\end{frame}

\begin{frame}{Arquivos}
  \huge
  \textbf{As funções \texttt{fprintf} e \texttt{fscanf} convertem os dados para caracteres, o que faz com que os arquivos sejam maiores e as operações levem mais tempo para serem executadas}
\end{frame}

\begin{frame}[fragile]{Arquivos}
  \huge
  \textbf{Movendo-se dentro do arquivo: \texttt{fseek}}
  \vfill
  \large
  \begin{lstlisting}[language=C]
int fseek(FILE *fp, long numbytes, int origem)
  \end{lstlisting}
  \vfill
  \setlength{\leftmargini}{0pt}
  \begin{itemize}
      \item [] Retorna \texttt{\textbf{0}}, caso a movimentação seja bem sucedida, ou outro valor, caso contrário
      \item [] \texttt{\textbf{origem}} pode ser \texttt{\textbf{SEEK\_SET}}, \texttt{\textbf{SEEK\_CUR}} ou \texttt{\textbf{SEEK\_END}}
      \item [] \texttt{\textbf{numbytes}} pode ser negativo
  \end{itemize}
\end{frame}

\begin{frame}[fragile]{Arquivos}
  \huge
  \textbf{Voltando ao início do arquivo: \texttt{rewind}}
  \vfill
  \large
  \begin{lstlisting}[language=C]
void rewind(FILE *fp)
  \end{lstlisting}
  \vfill
  \setlength{\leftmargini}{0pt}
  \begin{itemize}
      \item []
  \end{itemize}
\end{frame}

\begin{frame}[fragile]{Arquivos}
  \huge
  \textbf{Excluindo um arquivo: \texttt{remove}}
  \vfill
  \large
  \begin{lstlisting}[language=C]
int remove(char *nome_do_arquivo)
  \end{lstlisting}
  \vfill
  \setlength{\leftmargini}{0pt}
  \begin{itemize}
      \item [] Recebe o \textbf{caminho para o arquivo} e não o ponteiro para FILE
      \item [] Retorna \texttt{\textbf{0}}, caso o arquivo seja removido com sucesso, ou outro valor, caso contrário
  \end{itemize}
\end{frame}

\begin{frame}[fragile]{Arquivos}
  \huge
  \textbf{Erro ao acessar o arquivo: \texttt{ferror} e \texttt{perror}}
  \vfill
  \large
  \begin{lstlisting}[language=C]
int ferror(FILE *fp)
void perror(char *str)
  \end{lstlisting}
  \vfill
  \setlength{\leftmargini}{0pt}
  \begin{itemize}
      \item [] \texttt{ferror} detecta se a última operação realizada no arquivo produziu algum erro e retorna \texttt{\textbf{0}}, no caso negativo
      \item [] \texttt{perror} imprime na tela a string \textbf{\texttt{str}} seguida da mensagem de erro do sistema
  \end{itemize}
\end{frame}
\end{document}
